% ---- score2 standard preamble (from: python3 score2_tools.py preamble) ----
\documentclass[10pt]{extarticle}
\usepackage[a4paper,margin=1.5cm]{geometry}
\usepackage[T1]{fontenc}
\usepackage{amsmath,amssymb}
\usepackage{graphicx}
\usepackage[export]{adjustbox}
\usepackage{multicol}
\usepackage{enumitem}
\usepackage{booktabs,array,tabularx}
\setlength{\columnsep}{0.7cm}
\setlength{\parindent}{0pt}
\setlength{\parskip}{3pt}
\pagestyle{empty}
\makeatletter
% figure registry (figdefs.tex is generated by the extract step)
\newcommand{\qfigdef}[4]{\@namedef{qf@file@#1}{#2}\@namedef{qf@w@#1}{#3}\@namedef{qf@wide@#1}{#4}}
\newcommand{\qf@check}[1]{\@ifundefined{qf@file@#1}{\PackageError{score2}{Unknown figure id #1}{}}{}}
\newcommand{\qf@img}[4][]{\includegraphics[width=#2\dimexpr\@nameuse{qf@w@#4}pt\relax,max width=#3,max height=0.45\textheight,#1]{\@nameuse{qf@file@#4}}}
% Full-width questions: finalize wraps every question that contains a wide figure in \qfullbegin...\qfullend,
% so the WHOLE question (stem, figures, options) is set across the page in reading order, then the two columns
% resume. The question's head - its number, stem and first figure (\qfullsplit marks the end of the head) -
% is kept on one page, so a number is never left at a page foot with its figure overleaf; the rest may break
% across pages normally, so no big gaps are left. Without a \qfullsplit, the whole question is kept together
% if it is under 45% of a page. Never write these yourself.
\newif\ifqf@full
\newif\ifqf@boxed
\newsavebox{\qf@box}
\newcommand{\qf@open}{\global\qf@fulltrue\par\global\qf@boxedtrue
  \setbox\qf@box\vbox\bgroup\hsize\textwidth\linewidth\textwidth}
\newcommand{\qf@emit}[1]{\par\egroup\global\qf@boxedfalse
  \ifdim\dimexpr\ht\qf@box+\dp\qf@box\relax<#1\textheight\noindent\usebox{\qf@box}\par\else\unvbox\qf@box\fi}
\newcommand{\qfullsplit}{\ifqf@boxed\qf@emit{0.85}\fi}
\newcommand{\qf@close}{\ifqf@boxed\qf@emit{0.45}\fi\global\qf@fullfalse}
\newcommand{\qfullbegin}{\end{multicols}\qf@open}
\newcommand{\qfullend}{\qf@close\begin{multicols}{2}}
\newcommand{\qfullstart}{\qf@open}
\newcommand{\qfulldone}{\qf@close}
% \qfig[scale]{id}: block figure, centred, directly where it is written (size set by the tool).
% Fallback if a wide figure is not inside a full-width question: break out of the columns in place - no floats.
\newcommand{\qfig}[2][1]{\qf@check{#2}%
  \ifqf@full\qf@blk{#1}{#2}\else
  \ifnum\@nameuse{qf@wide@#2}=1 \expandafter\@firstoftwo\else\expandafter\@secondoftwo\fi
  {\end{multicols}\begin{center}\qf@img{#1}{\textwidth}{#2}\end{center}\begin{multicols}{2}}%
  {\qf@blk{#1}{#2}}\fi}
\newcommand{\qf@blk}[2]{\par\smallskip{\centering\qf@img{#1}{\linewidth}{#2}\par}\smallskip}
% \qfigi[scale]{id}: inline figure (inside option lists / tables / minipages); never breaks columns.
\newcommand{\qfigi}[2][1]{\qf@check{#2}\qf@img[valign=c]{#1}{\linewidth}{#2}}
\makeatother
% structure macros - lint reads these to check completeness and figure placement
\newcommand{\papersec}[1]{\par\bigskip{\centering\Large\bfseries PAPER-#1\par}\medskip}
\newcommand{\partsec}[2]{\par\medskip\noindent{\large\bfseries PART-#1: #2}\par\smallskip}
\newcommand{\sect}[1]{\par\smallskip\noindent\textbf{#1}\par}
\newcommand{\qn}[1]{\par\smallskip\noindent\textbf{#1.}~}
\newcommand{\paraq}[2]{\par\smallskip\noindent\textbf{Paragraph for Questions #1 and #2}\par}
\newenvironment{qoptions}{\begin{itemize}[leftmargin=1.8em,itemsep=1pt,topsep=2pt,label={}]}{\end{itemize}}
\newcommand{\opt}[2]{\item[(#1)] #2}
% figure registry (generated by score2_tools.py finalize)
\qfigdef{p1_phy_s1_q01_1}{Score-II - Test 17 (13 Feb 2024)-img/p1_phy_s1_q01_1.png}{169.2}{0}
\qfigdef{p1_phy_s1_q02_1}{Score-II - Test 17 (13 Feb 2024)-img/p1_phy_s1_q02_1.png}{109.4}{0}
\qfigdef{p1_phy_s1_q04_1}{Score-II - Test 17 (13 Feb 2024)-img/p1_phy_s1_q04_1.png}{131.3}{0}
\qfigdef{p1_phy_s1_q13_1}{Score-II - Test 17 (13 Feb 2024)-img/p1_phy_s1_q13_1.png}{199.7}{0}
\qfigdef{p1_phy_s1_q18_1}{Score-II - Test 17 (13 Feb 2024)-img/p1_phy_s1_q18_1.png}{154.1}{0}
\qfigdef{p1_phy_s1_q18_A}{Score-II - Test 17 (13 Feb 2024)-img/p1_phy_s1_q18_A.png}{153.8}{0}
\qfigdef{p1_phy_s1_q18_2}{Score-II - Test 17 (13 Feb 2024)-img/p1_phy_s1_q18_2.png}{154.3}{0}
\qfigdef{p1_phy_s1_q18_3}{Score-II - Test 17 (13 Feb 2024)-img/p1_phy_s1_q18_3.png}{154.3}{0}
\qfigdef{p1_phy_s1_q18_4}{Score-II - Test 17 (13 Feb 2024)-img/p1_phy_s1_q18_4.png}{154.3}{0}
\qfigdef{p1_phy_s1_q19_1}{Score-II - Test 17 (13 Feb 2024)-img/p1_phy_s1_q19_1.png}{79.0}{0}
\qfigdef{p1_phy_s1_q19_2}{Score-II - Test 17 (13 Feb 2024)-img/p1_phy_s1_q19_2.png}{79.4}{0}
\qfigdef{p1_phy_s1_q19_3}{Score-II - Test 17 (13 Feb 2024)-img/p1_phy_s1_q19_3.png}{79.0}{0}
\qfigdef{p1_phy_s1_q19_4}{Score-II - Test 17 (13 Feb 2024)-img/p1_phy_s1_q19_4.png}{79.7}{0}
\qfigdef{p1_phy_s1_q20_1}{Score-II - Test 17 (13 Feb 2024)-img/p1_phy_s1_q20_1.png}{124.3}{0}
\qfigdef{p1_phy_s1_q02_2}{Score-II - Test 17 (13 Feb 2024)-img/p1_phy_s1_q02_2.png}{133.2}{0}
\qfigdef{p1_phy_s1_q07_1}{Score-II - Test 17 (13 Feb 2024)-img/p1_phy_s1_q07_1.png}{154.1}{0}
\qfigdef{p1_chem_s1_q01_1}{Score-II - Test 17 (13 Feb 2024)-img/p1_chem_s1_q01_1.png}{113.0}{0}
\qfigdef{p1_chem_s1_q15_1}{Score-II - Test 17 (13 Feb 2024)-img/p1_chem_s1_q15_1.png}{154.6}{0}
\qfigdef{p1_chem_s1_q15_2}{Score-II - Test 17 (13 Feb 2024)-img/p1_chem_s1_q15_2.png}{71.8}{0}
\qfigdef{p1_chem_s1_q15_3}{Score-II - Test 17 (13 Feb 2024)-img/p1_chem_s1_q15_3.png}{79.2}{0}
\qfigdef{p1_chem_s1_q15_5}{Score-II - Test 17 (13 Feb 2024)-img/p1_chem_s1_q15_5.png}{91.4}{0}
\qfigdef{p1_chem_s1_q15_4}{Score-II - Test 17 (13 Feb 2024)-img/p1_chem_s1_q15_4.png}{91.9}{0}
\qfigdef{p1_chem_s1_q16_1}{Score-II - Test 17 (13 Feb 2024)-img/p1_chem_s1_q16_1.png}{183.6}{0}
\qfigdef{p1_chem_s1_q16_2}{Score-II - Test 17 (13 Feb 2024)-img/p1_chem_s1_q16_2.png}{77.0}{0}
\qfigdef{p1_chem_s1_q16_3}{Score-II - Test 17 (13 Feb 2024)-img/p1_chem_s1_q16_3.png}{77.3}{0}
\qfigdef{p1_chem_s1_q16_4}{Score-II - Test 17 (13 Feb 2024)-img/p1_chem_s1_q16_4.png}{78.5}{0}
\qfigdef{p1_chem_s1_q17_1}{Score-II - Test 17 (13 Feb 2024)-img/p1_chem_s1_q17_1.png}{199.0}{0}
\qfigdef{p1_chem_s1_q18_1}{Score-II - Test 17 (13 Feb 2024)-img/p1_chem_s1_q18_1.png}{169.4}{0}
\qfigdef{p1_chem_s1_q18_2}{Score-II - Test 17 (13 Feb 2024)-img/p1_chem_s1_q18_2.png}{30.0}{0}
\qfigdef{p1_chem_s1_q18_B}{Score-II - Test 17 (13 Feb 2024)-img/p1_chem_s1_q18_B.png}{30.7}{0}
\qfigdef{p1_chem_s1_q18_3}{Score-II - Test 17 (13 Feb 2024)-img/p1_chem_s1_q18_3.png}{64.3}{0}
\qfigdef{p1_chem_s1_q18_4}{Score-II - Test 17 (13 Feb 2024)-img/p1_chem_s1_q18_4.png}{56.9}{0}
\qfigdef{p1_chem_s1_q19_1}{Score-II - Test 17 (13 Feb 2024)-img/p1_chem_s1_q19_1.png}{206.9}{0}
\qfigdef{p1_chem_s1_q19_2}{Score-II - Test 17 (13 Feb 2024)-img/p1_chem_s1_q19_2.png}{216.2}{0}
\qfigdef{p1_chem_s1_q20_1}{Score-II - Test 17 (13 Feb 2024)-img/p1_chem_s1_q20_1.png}{53.0}{0}
\qfigdef{p1_chem_s1_q20_2}{Score-II - Test 17 (13 Feb 2024)-img/p1_chem_s1_q20_2.png}{52.8}{0}
\qfigdef{p1_chem_s1_q20_3}{Score-II - Test 17 (13 Feb 2024)-img/p1_chem_s1_q20_3.png}{72.5}{0}
\qfigdef{p1_chem_s2_q08_1}{Score-II - Test 17 (13 Feb 2024)-img/p1_chem_s2_q08_1.png}{154.1}{0}
\qfigdef{p1_chem_s2_q09_1}{Score-II - Test 17 (13 Feb 2024)-img/p1_chem_s2_q09_1.png}{131.3}{0}
\qfigdef{p1_chem_s2_q10_1}{Score-II - Test 17 (13 Feb 2024)-img/p1_chem_s2_q10_1.png}{177.1}{0}
\qfigdef{p1_math_s1_q04_1}{Score-II - Test 17 (13 Feb 2024)-img/p1_math_s1_q04_1.png}{50.9}{0}
\raggedcolumns
\makeatletter\@beginparpenalty=10000\makeatother
\begin{document}
\begin{center}{\Large\bfseries JEE (Main) Score-II --- Test 17 (13 Feb 2024)}\end{center}
\begin{multicols}{2}

\papersec{1}
\partsec{1}{PHYSICS}
\sect{Section-I}

\qn{1} In the circuit shown, the cell has emf = 10V and internal resistance = 1$\Omega$ :-
\qfig{p1_phy_s1_q01_1}
\begin{qoptions}
\opt{A}{The current through the 3$\Omega$ resistor is 2A}
\opt{B}{The corrent through the 3$\Omega$ resistor is 0.5A}
\opt{C}{The current through the 4$\Omega$ resistor is 0.5A}
\opt{D}{The current through the 4$\Omega$ resistor is 0.25A}
\end{qoptions}

\qn{2} Find the charge on the capacitor C in the following circuit :-
\qfig{p1_phy_s1_q02_1}
\begin{qoptions}
\opt{A}{12 $\mu$C}
\opt{B}{14 $\mu$C}
\opt{C}{20 $\mu$C}
\opt{D}{18 $\mu$C}
\end{qoptions}

\qn{3} A bar magnet when placed at an angle of $30^\circ$ to the direction of magnetic field induction of $5\times10^{-2}$T, experiences a moment of couple $25\times10^{-6}$ N-m. If the length of the magnet is 5 cm its pole strength is:-
\begin{qoptions}
\opt{A}{$2\times10^{-2}$ A-m}
\opt{B}{$5\times10^{-2}$ A-m}
\opt{C}{2 A-m}
\opt{D}{5 A-m}
\end{qoptions}

\qn{4} A conducting square loop of side L and resistance R moves in its plane with a uniform velocity v perpendicular to one of its sides. A magnetic induction B constant in time and space, pointing perpendicular and into the plane of the loop exists everywhere. The current induced in the loop is :
\qfig{p1_phy_s1_q04_1}
\begin{qoptions}
\opt{A}{$\dfrac{BLv}R$ clockwise}
\opt{B}{$\dfrac{BLv}R$ anticlockwise}
\opt{C}{$\dfrac{2BLv}R$ anticlockwise}
\opt{D}{Zero}
\end{qoptions}

\qn{5} In an ac circuit, V and I are given by\\
$V=100\sin(100\,t)$ volts,\\
$I=100\sin\left(100t+\dfrac\pi3\right)$mA\\
The power dissipated in circuit is :
\begin{qoptions}
\opt{A}{$10^4$ watt}
\opt{B}{10 watt}
\opt{C}{2.5 watt}
\opt{D}{5 watt}
\end{qoptions}

\qn{6} For a plane electromagnetic wave propagating in the $+Z$ direction, which one of the following combination gives the correct possible direction for $\vec E$ and $\vec B$ field respectively ?
\begin{qoptions}
\opt{A}{$\left(\hat i+2\hat j\right)$ and $\left(2\hat i-\hat j\right)$}
\opt{B}{$\left(-2\hat i-3\hat j\right)$ and $\left(3\hat i-2\hat j\right)$}
\opt{C}{$\left(2\hat i+3\hat j\right)$ and $\left(\hat i+2\hat j\right)$}
\opt{D}{$\left(3\hat i+4\hat j\right)$ and $\left(4\hat i-3\hat j\right)$}
\end{qoptions}

\qn{7} The potential energy of a particle varies with distance x as $V=\dfrac{A\sqrt x}{x+B}$. The dimensions of AB are :-
\begin{qoptions}
\opt{A}{$ML^{5/2}T^{-2}$}
\opt{B}{$M^1L^2T^{-2}$}
\opt{C}{$M^{3/2}L^{3/2}T^{-2}$}
\opt{D}{$M^1L^{7/2}T^{-2}$}
\end{qoptions}

\qn{8} The equation of motion of a projectile are x = 36t and 2y = 96t $-$ 9.8t$^2$. Angle of projectionis :-
\begin{qoptions}
\opt{A}{$\sin^{-1}\left(\dfrac45\right)$}
\opt{B}{$\sin^{-1}\left(\dfrac35\right)$}
\opt{C}{$\sin^{-1}\left(\dfrac43\right)$}
\opt{D}{$\sin^{-1}\left(\dfrac34\right)$}
\end{qoptions}

\qn{9} A body sliding on a smooth inclined plane requires 4 sec to reach the bottom after starting from rest at the top. How much time does it take to cover one fourth the distance starting from the top :-
\begin{qoptions}
\opt{A}{1 sec}
\opt{B}{2 sec}
\opt{C}{0.4 sec}
\opt{D}{1.6 sec}
\end{qoptions}

\qn{10} A rod of length L has one end on horizontal floor. It falls while rotating about contact point without slipping. Initially it is inclined at an angle $\alpha$ from horizontal. Its angular speed when it is at horizontal floor is :-
\begin{qoptions}
\opt{A}{$\sqrt{\dfrac{3g\sin\alpha}L}$}
\opt{B}{$\sqrt{\dfrac{2L}{3g\sin\alpha}}$}
\opt{C}{$\sqrt{\dfrac{6g\sin\alpha}L}$}
\opt{D}{$\sqrt{\dfrac L{g\sin\alpha}}$}
\end{qoptions}

\qn{11} The orbital velocity of an artificial satellite in a circular orbit very close to earth is V. The velocity of a geo-stationary satellite orbiting in circular orbit at an altitude of 3R from earth's surface will be :-
\begin{qoptions}
\opt{A}{$\dfrac V{\sqrt7}$}
\opt{B}{$\dfrac V{\sqrt6}$}
\opt{C}{$\dfrac V2$}
\opt{D}{$\dfrac V{\sqrt2}$}
\end{qoptions}

\qn{12} One end of a uniform wire of length L and of weight W is attached rigidly to a point in the roof and a weight $W_1$ is suspended from its lower end. If S is the area of cross-section of the wire, the stress in the wire at a height 3L/4 from its lower end is :
\begin{qoptions}
\opt{A}{$\dfrac{W_1}S$}
\opt{B}{$\dfrac{W_1+(W/4)}S$}
\opt{C}{$\dfrac{W_1+(3W/4)}S$}
\opt{D}{$\dfrac{W_1+W}S$}
\end{qoptions}

\qn{13} Consider slabs of three media A, B and C arranged as shown in figure. R.I. of A is 1.5 and that of C is 1.4. If the number of waves in A is equal to the number of waves in the combination of B and C then refractive index of B is :-
\qfig{p1_phy_s1_q13_1}
\begin{qoptions}
\opt{A}{1.4}
\opt{B}{1.5}
\opt{C}{1.6}
\opt{D}{1.7}
\end{qoptions}

\qn{14} White light is used to illuminate the two slits in a Young's double slit experiment. The separation between slits is b and the screen is at a distance [d $\gg$ b] from the slits. At a point on the screen directly infront of one of the slits, certain wavelengths are missing. Some of these missing wavelengths are-\\
\relax[a] $\lambda=b^2/d$ \hfill [b] $\lambda=2b^2/d$\\
\relax[c] $\lambda=b^2/3\,d$ \hfill [d] $\lambda=2b^2/3d$
\begin{qoptions}
\opt{A}{a, b}
\opt{B}{b, c}
\opt{C}{a, c}
\opt{D}{a, d}
\end{qoptions}

\qn{15} Cut-off potentials for a metal surface for light of wavelength $\lambda_1$, $\lambda_2$ and $\lambda_3$ are $V_1$, $V_2$ and $V_3$ respectively. If $V_1$, $V_2$ and $V_3$ are in arithmetic progression, then $\lambda_1$, $\lambda_2$ and $\lambda_3$ are in:-
\begin{qoptions}
\opt{A}{arithmetic progression}
\opt{B}{geometric progression}
\opt{C}{harmonic progression}
\opt{D}{None of these}
\end{qoptions}

\qn{16} The binding energies per nucleon for a deutron and an $\alpha$ - particle are $E_1$ and $E_2$ respectively. What will be the energy E released in the following reaction ?\\
${}_1\text{H}^2+{}_1\text{H}^2\rightarrow{}_2\text{H}^4+E$
\begin{qoptions}
\opt{A}{$2(E_1+E_2)$}
\opt{B}{$2(E_2-E_1)$}
\opt{C}{$4(E_1+E_2)$}
\opt{D}{$4(E_2-E_1)$}
\end{qoptions}

\qn{17} The minimum light intensity that can be perceived by human eye is about $10^{-10}$ Wm$^{-2}$. If the area of the pupil of our eye is 0.8 mm$^2$ and the average wavelength of visible light is 5000 \AA{} then the number of visible light photons entering per second in our pupil will be :-
\begin{qoptions}
\opt{A}{$10^2$}
\opt{B}{$2\times10^2$}
\opt{C}{$10^4$}
\opt{D}{$2\times10^4$}
\end{qoptions}

\qn{18} The logic circuit shown has the input waveforms `A' and `B' as shown. Pick out the \textbf{CORRECT} output waveform :-
\qfig{p1_phy_s1_q18_1}
\begin{qoptions}
\opt{A}{\qfigi{p1_phy_s1_q18_A}}
\opt{B}{\qfigi{p1_phy_s1_q18_2}}
\opt{C}{\qfigi{p1_phy_s1_q18_3}}
\opt{D}{\qfigi{p1_phy_s1_q18_4}}
\end{qoptions}

\qn{19} In the following figures heat is absorbed by the gas :-
\begin{qoptions}
\opt{A}{\qfigi{p1_phy_s1_q19_1}}
\opt{B}{\qfigi{p1_phy_s1_q19_2}}
\opt{C}{\qfigi{p1_phy_s1_q19_3}}
\opt{D}{\qfigi{p1_phy_s1_q19_4}}
\end{qoptions}

\qn{20} Six identical rods are arranged as shown in the figure. The temperature of the junction B will be :-
\qfig{p1_phy_s1_q20_1}
\begin{qoptions}
\opt{A}{$\dfrac{240^\circ}3$C}
\opt{B}{$240^\circ$C}
\opt{C}{$140^\circ$C}
\opt{D}{$\dfrac{140^\circ}3$C}
\end{qoptions}

\sect{Section-II}

\qn{1} The electric potential in some region is expressed by V = 6x $-$ 8xy$^2$ $-$ 8y + 6yz $-$ 4z$^2$ volt. The magnitude of force acting on a charge of 2 C situated at the origin will be :-

\qn{2} In a meter bridge experiment S is a standard resistance. R is a resistance wire. It Is found that balancing length is $\ell$ = 25 cm. If R is replaced by a wire of half length and half diameter that of R of same material, then the balancing distance $\ell'$ (in cm) will now be\underline{\hspace{1cm}}.
\qfig{p1_phy_s1_q02_2}

\qn{3} A particle moves in the x-y plane under the action of a force $\vec F$ such that the value of its linear momentum $\vec p$ at any time t is $p_x=2\cos t$ and $p_y=2\sin t$. The angle $\theta$ between $\vec F$ and $\vec p$ at the time t is (in degree) :-

\qn{4} The potential energy between two atoms in a molecule is given by $U_{(x)}=\dfrac{16}{x^{12}}-\dfrac1{2x^6}$, where x is the distance between the atoms. The system is in equilibrium when x is equal to\dots\dots\dots

\qn{5} An object is placed on the axis of a convex lens. Its image is formed 80 cm from the object. The magnification is 3. The focal length of the lens (in cm) is :-

\qn{6} Velocity of sound in air is 330 m/s. For a particular sound in air, a path difference of 40 cm is equivalent to a phase difference of $1.6\pi$. The frequency of the wave (in Hz) is :-

\qn{7} In the given circuit, the diode has a forward resistance of 8 $\Omega$ and infinite backward resistance. The value of R is :-
\qfig{p1_phy_s1_q07_1}

\qn{8} A nucleus of an element ${}_{84}X^{202}$ emits an $\alpha-$ particle first, $\beta-$particle next and then a gamma photon. The final nucleus formed has an atomic number.

\qn{9} Two ideal monoatomic \& diatomic gases are mixed with one another to form an ideal gas mixture. The equation of the adiabatic process of the mixture is $PV^\gamma$ = constant, where $\gamma$ = 11/7. If $n_1$ \& $n_2$ are the number of moles of the monoatomic \& diatomic gases in the mixture respectively, find the ratio $n_1/n_2$ :-

\qn{10} A black body at a temperature of 1640 K has the wavelength corresponding to maximum emission equal to 1.75$\mu$m. Assuming the moon to be a perfectly black body, the temperature of the moon, if the wavelength corresponding to maximum emission is 14.35$\mu$m is ( in kelvin)

\partsec{2}{CHEMISTRY}
\sect{Section-I}

\qn{1} The preparation of SO$_3$(g) by reaction\\
SO$_2$(g) + $\dfrac12$ O$_2$(g) $\rightleftharpoons$ SO$_3$(g)
\qfig{p1_chem_s1_q01_1}
is an exothermic reaction. If the preparation follows the following temperature-pressure relationship for its \% yield, then for temperatures $T_1$, $T_2$ and $T_3$. The correct option is :-
\begin{qoptions}
\opt{A}{$T_3>T_2>T_1$}
\opt{B}{$T_1>T_2>T_3$}
\opt{C}{$T_1=T_2=T_3$}
\opt{D}{Nothing could be predicted about temperature through given information.}
\end{qoptions}

\qn{2} In the following first order reactions\\
(A) $\xrightarrow{K_1}$ Product, (B) $\xrightarrow{K_2}$ Product , find the ratio $K_1/K_2$ if 90\% of (A) has been reacted in time `t' while 99\% of (B) has been reacted in time 2t :-
\begin{qoptions}
\opt{A}{1}
\opt{B}{2}
\opt{C}{1/2}
\opt{D}{$2\left(\dfrac{1-2\log3}{2-\log11-2\log3}\right)$}
\end{qoptions}

\qn{3} The electrolysis of acetate solution produces ethane according to reaction\\
$2\text{CH}_3\text{COO}^-\rightarrow\text{C}_2\text{H}_6(g)+2\text{CO}_2(g)+2e^-$\\
The current efficiency of the process is 80\%. What volume of gases would be produced at 27$^\circ$C and 740 torr, if the current of 0.5 amp is passed through the solution for 96.5 min ?
\begin{qoptions}
\opt{A}{6.0 L}
\opt{B}{0.60 L}
\opt{C}{1.365 L}
\opt{D}{0.91 L}
\end{qoptions}

\qn{4} 20 ml, 0.1M HA $\left[(K_b)_{A^-}=10^{-10}\right]$ is titrated against 0.2M, NaOH. Select the correct statement :-\\
\relax[Given : log2 = 0.3, log3 = 0.48]
\begin{qoptions}
\opt{A}{Initially pH of HA solution is 11.5}
\opt{B}{At half equivalence point pH of solution is 4 and at equivalence point pH is 9.09}
\opt{C}{Phenolphthalein is suitable indicator for this titration but methyl orange is not suitable}
\opt{D}{Phenolphthalein \& methyl orange both are suitable indicator for end point determination.}
\end{qoptions}

\qn{5} If the elevation in boiling point of a solution of non volatile, non-electrolytic and non-associating solute in solvent ($K_b$ = xK kg.mol$^{-1}$) is y K, then the depression in freezing point of solution of same concentration would be :\\
($K_f$ of the solvent = zK kg.mol$^{-1}$) -
\begin{qoptions}
\opt{A}{$\dfrac{2xz}y$}
\opt{B}{$\dfrac{yz}x$}
\opt{C}{$\dfrac{xz}y$}
\opt{D}{$\dfrac{yz}{2x}$}
\end{qoptions}

\qn{6} Enthalpy of neutralization is defined as the enthalpy change when 1gm equivalent of acid/base is completely neutralized by base/acid in dilute solution. For strong acid and strong base neutralization net chemical change is\\
H$^+$(aq) + OH$^-$(aq) $\rightarrow$ H$_2$O($\ell$) ;\\
$\Delta_rH^\circ=-55.84$kJ/mol \ $\Delta H^\circ_{\text{ionization}}$ of aqueous solution of strong acid and strong base is zero. When a dilute solution of a weak acid or based is neutralized, the enthalpy of neutralization is some what less because of the absorption of heat in the ionization of the weak acid or base, for weak acid/base\\
$\Delta H^\circ_{\text{neutralization}}=\Delta H^\circ_{\text{ionization}}+\Delta_rH^\circ(\text{H}^++\text{OH}^-\rightarrow\text{H}_2\text{O})$\\
Under the same conditions how many mL of 0.1 M NaOH and 0.05 M H$_2$A (strong diprotic acid) solution should be mixed for a total volume of 100mL produce the highest rise in temperature :-
\begin{qoptions}
\opt{A}{25 : 75}
\opt{B}{50 : 50}
\opt{C}{75 : 25}
\opt{D}{66.66 : 33.33}
\end{qoptions}

\qn{7} The following reaction is at equilibrium at 298K.\\
2NO(0.00001 bar) + Cl$_2$(g, 0.01 bar)\\
\hspace*{1.2cm}$\rightleftharpoons$ 2NOCl (g, 0.01 bar)\\
$\Delta G^o$ for the reaction is :-
\begin{qoptions}
\opt{A}{$-45.65$ kJ}
\opt{B}{$-28.53$ kJ}
\opt{C}{$-22.83$ kJ}
\opt{D}{$-57.06$ kJ}
\end{qoptions}

\qn{8} In which change bond order and no. of unpaired electron(s) both are increased ?
\begin{qoptions}
\opt{A}{N$_2$ $\rightarrow$ N$_2^+$}
\opt{B}{O$_2$ $\rightarrow$ O$_2^+$}
\opt{C}{O$_2^-$ $\rightarrow$ O$_2^{2-}$}
\opt{D}{O$_2^-$ $\rightarrow$ O$_2$}
\end{qoptions}

\qn{9} Which of the following statements is true on the basis of back bonding ?
\begin{qoptions}
\opt{A}{Si$-$O bond is stronger than C$-$O bond}
\opt{B}{Dimethyl ether acts a better lewis base than disilyl ether (SiH$_3-$O$-$SiH$_3$)}
\opt{C}{(CH$_3$)$_3$C$-$O$-$H is less acidic than (CH$_3$)$_3$Si$-$O$-$H}
\opt{D}{All of these}
\end{qoptions}

\qn{10} Which of the following having sp hybridization ?
\begin{qoptions}
\opt{A}{I$_3^-$}
\opt{B}{NO$_3^-$}
\opt{C}{PCl$_5$}
\opt{D}{CO$_2$}
\end{qoptions}

\qn{11} Consider the complex ion, $[\text{Cr(NH}_3)_2\text{Cl}_2(\text{C}_2\text{O}_4)]^-$. What is the :-\\
(i) Oxidation number of the metal ion ?\\
(ii) Co-ordination number of the metal ion ?\\
(iii) Charge on the complex if all ligands were chloride ions ?\par\nobreak
\begin{center}\small
\begin{tabular}{|c|c|c|c|}
\hline
 & (i) & (ii) & (iii) \\
\hline
(A) & $-1$ & 6 & $-3$ \\
\hline
(B) & $+3$ & 6 & $-3$ \\
\hline
(C) & $+3$ & 6 & $-1$ \\
\hline
(D) & $-1$ & 5 & $-3$ \\
\hline
\end{tabular}
\end{center}

\qn{12} In the following complex ions order of paramagnetism is :\\
P : $[\text{FeF}_6]^{3-}$ \hspace{1cm} Q : $[\text{CoF}_6]^{3-}$\\
R : $[\text{V(H}_2\text{O)}_6]^{3+}$ \hspace{0.5cm} S : $[\text{Ti(H}_2\text{O)}_6]^{3+}$
\begin{qoptions}
\opt{A}{P $>$ Q $>$ R $>$ S}
\opt{B}{Q $>$ P $>$ R $>$ S}
\opt{C}{P = Q = R = S}
\opt{D}{P $>$ R $>$ Q $>$ S}
\end{qoptions}

\qn{13} The type of isomerism present in $[\text{Cr(NH}_3)_5\,(\text{NO}_2)]\,\text{Cl}_2$ is :
\begin{qoptions}
\opt{A}{optical}
\opt{B}{linkage}
\opt{C}{hydrate}
\opt{D}{polymerisation}
\end{qoptions}

\qn{14} The correct order of ionic radii of $\text{Y}^{3+}$, $\text{La}^{3+}$, $\text{Eu}^{3+}$ and $\text{Lu}^{3+}$ is :-
\begin{qoptions}
\opt{A}{$\text{Y}^{3+}<\text{La}^{3+}<\text{Eu}^{3+}<\text{Lu}^{3+}$}
\opt{B}{$\text{Y}^{3+}<\text{Lu}^{3+}<\text{Eu}^{3+}<\text{La}^{3+}$}
\opt{C}{$\text{Lu}^{3+}<\text{Eu}^{3+}<\text{La}^{3+}<\text{Y}^{3+}$}
\opt{D}{$\text{La}^{3+}<\text{Eu}^{3+}<\text{Lu}^{3+}<\text{Y}^{3+}$}
\end{qoptions}

\qn{15} \qfig{p1_chem_s1_q15_1}
(A) , Product (A) is :
\begin{qoptions}
\opt{A}{\qfigi{p1_chem_s1_q15_2}}
\opt{B}{\qfigi{p1_chem_s1_q15_3}}
\opt{C}{\qfigi{p1_chem_s1_q15_5}}
\opt{D}{\qfigi{p1_chem_s1_q15_4}}
\end{qoptions}

\qn{16} Identify C in the reactions :
\qfig{p1_chem_s1_q16_1}
\begin{qoptions}
\opt{A}{\qfigi{p1_chem_s1_q16_2}}
\opt{B}{\qfigi{p1_chem_s1_q16_3}}
\opt{C}{\qfigi{p1_chem_s1_q16_4}}
\opt{D}{None of these}
\end{qoptions}

\qn{17} Compound (X) reacts with NaNO$_2$ \& H$_2$SO$_4$ \& then reacts with ortho cresol to form sharp colour compound, Compound X may be :-
\qfig{p1_chem_s1_q17_1}
\begin{qoptions}
\opt{A}{a and b}
\opt{B}{a, b and c}
\opt{C}{b and d}
\opt{D}{d only}
\end{qoptions}

\qn{18} \qfig{p1_chem_s1_q18_1}
Product (C) is :-
\begin{qoptions}
\opt{A}{\qfigi{p1_chem_s1_q18_2}}
\opt{B}{\qfigi{p1_chem_s1_q18_B}}
\opt{C}{\qfigi{p1_chem_s1_q18_3}}
\opt{D}{\qfigi{p1_chem_s1_q18_4}}
\end{qoptions}

\qn{19} \qfig{p1_chem_s1_q19_1}
Correct statement about major product :-
\begin{qoptions}
\opt{A}{If follows S$_N$1 and racemic mixture of alcohol is obtained.}
\opt{B}{It follows S$_N$1 and racemic mixture of alkyl iodide is obtained.}
\opt{C}{It follows S$_N$2 and racemic mixture of alcohol is obtained.}
\opt{D}{If follows S$_N$2 and racemic mixture of alkyl iodide is obtained.}
\end{qoptions}

\qn{20} \qfig{p1_chem_s1_q19_2}
Structure of (A) is :-
\begin{qoptions}
\opt{A}{\qfigi{p1_chem_s1_q20_1}}
\opt{B}{\qfigi{p1_chem_s1_q20_2}}
\opt{C}{\qfigi{p1_chem_s1_q20_3}}
\opt{D}{None of these}
\end{qoptions}

\sect{Section-II}

\qn{1} For a complex reaction A $\xrightarrow{k}$ products\\
$E_{a_1}$=180 kJ/mol; $E_{a_2}$ =80 kJ/mol; $E_{a_3}$ = 50 kJ/mol\\
Overall rate constant k is related to individual rate constant by the equation $k=\left(\dfrac{k_1.k_2}{k_3}\right)^{2/3}$\\
Activation energy (kJ/mol) for the overall reaction is :-

\qn{2} The standard electrode potential for the following reaction is +1.33V. What is the potential at pH = 2.0\\
$\text{Cr}_2\text{O}_7^{2-}$ (aq 1M) + 14H$^+$(aq) + 6e$^-$ $\rightarrow$ 2Cr$^{+3}$(aq,1M) + 7H$_2$O($\ell$)

\qn{3} A solution of isopropyl alcohol and propyl alcohol has a vapour pressure 200 mm of Hg if it has 25\% mole of isopropyl alcohol. Another solution of same components containing 25\% mole propyl alcohol has vapour pressure 300 mm of Hg. Then vapor pressure of isopropyl alcohol in mm of Hg is:-

\qn{4} What weight of AgCl will be formed when a solution containing 4.77 g of NaCl is added to a solution of 5.77 g of AgNO$_3$ ?

\qn{5} Define the oxidation states of Mn in product of the given reaction\\
$3\text{K}_2\text{MnO}_4+2\text{H}_2\text{O}+4\text{CO}_2\rightarrow2\text{X}+\text{Y}+4\text{KHCO}_3$\\
If the oxidation state of Mn in product X and Y are $n_1$ and $n_2$ respectively. Then find out the value of ($n_1+n_2$) -

\qn{6} Among the carbide how many carbide will produce methane gas on hydrolysis?\\
Al$_4$C$_3$, Be$_2$C, CaC$_2$, SrC$_2$, Mg$_2$C$_3$

\qn{7} How many of following orders are correct?\\
(a) Order of IE : O$^{-2}$ $>$ O$^+$ $>$ O $>$ O$^-$\\
(b) Order of EN : Zn $<$ Cd $<$ Hg\\
(c) Order of IE : B $>$ Tl $>$ Ga $>$ Al $>$ In\\
(d) Order of EA : S $>$ Se $>$ Te $>$ O\\
(e) Order of atomic radius : Ni $<$ Cu $<$ Zn\\
(f) Order of ionic radius : Na$^+$ $>$ Mg$^{+2}$ $>$ Li$^+$ $>$ Be$^{+2}$\\
(g) Order of IE$_2$ : N $<$ O $<$ F $<$ Ne

\qn{8} The total number of carboxylic groups in product P is/are
\qfig{p1_chem_s2_q08_1}

\qn{9} How many moles of HIO$_4$ are consumed by given compound when it react with HIO$_4$ :-
\qfig{p1_chem_s2_q09_1}

\qn{10} Sum of $\alpha$-hydrogen in major product of the reaction.
\qfig{p1_chem_s2_q10_1}
Sum of $\alpha$-hydrogen is A + B + C + D =

\partsec{3}{MATHEMATICS}
\sect{Section-I}

\qn{1} Give the correct order of initials T or F for following statements. Use T if statement is true and F if it is false.\\
Statement-1 : If the graphs of two linear equations in two variables are neither parallel nor identical, then there is a unique solution to the system.\\
Statement-2 : If the system of equations\\
ax + by = 0, cx + dy = 0 has a non-zero solution, then it has infinitely many solutions.\\
Statement-3 : The system of equations\\
x + y + z = 1, x = y, y = 1 + z is inconsistent.\\
Statement-4 : If two of the equations in a system of three linear equations are inconsistent, then the whole system is inconsistent.
\begin{qoptions}
\opt{A}{FFTT}
\opt{B}{TTFT}
\opt{C}{TTFF}
\opt{D}{TTTF}
\end{qoptions}

\qn{2} If $\alpha_1<\alpha_2<\alpha_3<\alpha_4<\alpha_5<\alpha_6$ , then the equation $(x-\alpha_1)(x-\alpha_3)(x-\alpha_5)+3(x-\alpha_2)(x-\alpha_4)(x-\alpha_6)=0$ has :-
\begin{qoptions}
\opt{A}{no real root in $(\alpha_5,\alpha_6)$}
\opt{B}{no real root in $(\alpha_1,\alpha_2)$}
\opt{C}{all roots are imaginary}
\opt{D}{no real root in $(-\infty,\alpha_1)$}
\end{qoptions}

\qn{3} The number of positive proper divisors of 1800, which are also divisible by 10 are
\begin{qoptions}
\opt{A}{34}
\opt{B}{17}
\opt{C}{16}
\opt{D}{18}
\end{qoptions}

\qn{4} Number of ways in which AAABBB can be placed in the squares of figure as shown so that no row remains empty is :-
\qfig{p1_math_s1_q04_1}
\begin{qoptions}
\opt{A}{2430}
\opt{B}{2160}
\opt{C}{1620}
\opt{D}{None of these}
\end{qoptions}

\qn{5} If two events A and B are such that $P(A^C)=0.3$, $P(B)=0.4$ and $P(A\cap B^C)=0.5$, then value of $P\left[\dfrac B{A\cup B^C}\right]$ is equal to :-
\begin{qoptions}
\opt{A}{$\dfrac14$}
\opt{B}{$\dfrac12$}
\opt{C}{$\dfrac34$}
\opt{D}{$\dfrac38$}
\end{qoptions}

\qn{6} In the expansion of $(167)^{131}$ the unit place digit is :$-$
\begin{qoptions}
\opt{A}{7}
\opt{B}{9}
\opt{C}{1}
\opt{D}{3}
\end{qoptions}

\qn{7} If $|z_1|=2$, $|z_2|=3$, $|z_3|=4$ and $|2z_1+3z_2+4z_3|=9$ , then value of $|8z_2z_3+27z_3z_1+64z_1z_2|$is equal to :-
\begin{qoptions}
\opt{A}{216}
\opt{B}{108}
\opt{C}{42}
\opt{D}{56}
\end{qoptions}

\qn{8} Let $f(x)=\max\{\sin^{-1}x,\cos^{-1}x\}$ then which of the following is not correct
\begin{qoptions}
\opt{A}{$f(x)$ is not differentiable at $x=\dfrac1{\sqrt2}$}
\opt{B}{$f(x)$ is continuous everywhere}
\opt{C}{Range of $f(x)$ is $\left[\dfrac\pi2,\pi\right]$}
\opt{D}{$f(x)$ is many one function}
\end{qoptions}

\qn{9} The solution of differential equation,\\
$e^{-x}(y+1)dy+y(\cos^2x-\sin2x)dx=0$\\
subject to condition y = 1 when x = 0, is :
\begin{qoptions}
\opt{A}{$(y+1)+e^x\cos^2x=3$}
\opt{B}{$\ell n(y+1)+e^x.\cos^2x=1+\ell n2$}
\opt{C}{$y+\ell ny=e^x\cos^2x$}
\opt{D}{$y+\ell ny+e^x.\cos^2x=2$}
\end{qoptions}

\qn{10} If $\displaystyle\int\dfrac{\cos x-\sin x+1-x}{e^x+\sin x+x}=\ln(f(x))+g(x)+C$\\
where C is the constant of integration and $f(x)$ is positive, then $f(x)+g(x)$ has the value equal to
\begin{qoptions}
\opt{A}{$e^x+\sin x+2x$}
\opt{B}{$e^x+\sin x$}
\opt{C}{$e^x-\sin x$}
\opt{D}{$e^x+\sin x+x$}
\end{qoptions}

\qn{11} If $\sin^{-1}\sqrt{\left(x^2+2x+1\right)}+\sec^{-1}\sqrt{\left(x^2+2x+1\right)}=\dfrac\pi2$,\\
$x\neq0$, then the value of\\
$2\sec^{-1}\left(\dfrac x2\right)+\sin^{-1}\left(\dfrac x2\right)$ is equal to
\begin{qoptions}
\opt{A}{$-\dfrac{3\pi}2$}
\opt{B}{$\dfrac{3\pi}2$}
\opt{C}{$-\dfrac\pi2$}
\opt{D}{$\dfrac\pi2$}
\end{qoptions}

\qn{12} Let a = Minimum $\{x^2+2x+3,\ x\in R\}$ and\\
$b=\displaystyle\lim_{\theta\to0}\dfrac{1-\cos\theta}{\theta^2}$. The value of $\displaystyle\sum_{r=0}^na^r.b^{n-r}$ is:-
\begin{qoptions}
\opt{A}{$\dfrac{2^{n+1}-1}{3.2^n}$}
\opt{B}{$\dfrac{2^{n+1}+1}{3.2^n}$}
\opt{C}{$\dfrac{4^{n+1}-1}{3.2^n}$}
\opt{D}{None of these}
\end{qoptions}

\qn{13} Two circles whose radii are equal to 4 and 8 intersect at right angles. The length of their common chord is-
\begin{qoptions}
\opt{A}{$\dfrac{16}{\sqrt5}$}
\opt{B}{8}
\opt{C}{$4\sqrt6$}
\opt{D}{$\dfrac{8\sqrt5}5$}
\end{qoptions}

\qn{14} A chord is drawn through the focus of the parabola $y^2=6x$ such that its distance from the vertex of this parabola is $\dfrac{\sqrt5}2$, then its slope can be :
\begin{qoptions}
\opt{A}{$\dfrac{\sqrt3}2$}
\opt{B}{$\dfrac{\sqrt5}2$}
\opt{C}{$\dfrac2{\sqrt3}$}
\opt{D}{$\dfrac2{\sqrt5}$}
\end{qoptions}

\qn{15} The hyperbola $\dfrac{x^2}{a^2}-\dfrac{y^2}{b^2}=1$ has its conjugate axis 5 and passes through the point $(2,-1)$. The length of a latus rectum is :
\begin{qoptions}
\opt{A}{$\dfrac34\sqrt{29}$}
\opt{B}{$\dfrac58\sqrt{29}$}
\opt{C}{$\sqrt{29}$}
\opt{D}{$\dfrac54\sqrt{29}$}
\end{qoptions}

\qn{16} A line makes angle $45^\circ$, $60^\circ$, and $60^\circ$ with the coordinate axis, the projection of line segments on line which joins point ($-$ 1,2,3) and ($-$ 1, 4, 0) are-
\begin{qoptions}
\opt{A}{3/2}
\opt{B}{1/3}
\opt{C}{1/2}
\opt{D}{2/3}
\end{qoptions}

\qn{17} If $|\vec c|^2=60$ and $\vec c\times\left(\hat i+2\hat j+5\hat k\right)=\vec0$, then a value of $\vec c.\left(-7\hat i+2\hat j+3\hat k\right)$ is :
\begin{qoptions}
\opt{A}{$4\sqrt2$}
\opt{B}{24}
\opt{C}{$12\sqrt2$}
\opt{D}{12}
\end{qoptions}

\qn{18} Let $S=\left\{x\in\left(-\dfrac\pi2,\dfrac\pi2\right):9^{1-\tan^2x}+9^{\tan^2x}=10\right\}$ and $\beta=\displaystyle\sum_{x\in S}\tan^2\left(\dfrac x3\right)$, then $\dfrac16(\beta-14)^2$ is equal to
\begin{qoptions}
\opt{A}{32}
\opt{B}{8}
\opt{C}{64}
\opt{D}{16}
\end{qoptions}

\qn{19} Let $\mu$ be the mean and $\sigma$ be the standard deviation of the distribution\par\nobreak
\begin{center}\small
\begin{tabular}{|c|c|c|c|c|c|c|}
\hline
$x_i$ & 0 & 1 & 2 & 3 & 4 & 5 \\
\hline
$f_i$ & $k+2$ & $2k$ & $k^2-1$ & $k^2-1$ & $k^2+1$ & $k-3$ \\
\hline
\end{tabular}
\end{center}
where $\sum f_i=62$. if [x] denotes the greatest integer $\leq$ x, then $[\mu^2+\sigma^2]$ is equal
\begin{qoptions}
\opt{A}{8}
\opt{B}{7}
\opt{C}{6}
\opt{D}{9}
\end{qoptions}

\qn{20} Let in a right angled triangle, the smallest angle be $\theta$. If a triangle formed by taking the reciprocal of its sides is also a right angled triangle, then $\sin\theta$ is equal to :
\begin{qoptions}
\opt{A}{$\dfrac{\sqrt5+1}4$}
\opt{B}{$\dfrac{\sqrt5-1}2$}
\opt{C}{$\dfrac{\sqrt2-1}2$}
\opt{D}{$\dfrac{\sqrt5-1}4$}
\end{qoptions}

\sect{Section-II}

\qn{1} The greatest integer less than or equal to the sum of first 100 terms of the sequence\\
$\dfrac13,\dfrac59,\dfrac{19}{27},\dfrac{65}{81},\dots$ is equal to

\qn{2} If $|A|_{3\times3}=9$, then value of $\left|\text{adj}\left(\dfrac A3\right)^{-1}\right|$ is $-$

\qn{3} The area enclosed by the curve $xy^2=4(2-x)$ and y-axis is $\lambda$ square unit, then the value of $8\cos(2014\,\lambda)$ is equal to

\qn{4} Number of values of x satisfying the equation\\
$\displaystyle\int_{-1}^x\left(8t^2+\dfrac{28}3t+4\right)dt=\dfrac{\left(\frac32\right)x+1}{\log_{(x+1)}\sqrt{x+1}}$, is :

\qn{5} Let $f(n)=\left[\dfrac13+\dfrac{3n}{100}\right]n$, where [n] denotes the greatest integer less than or equal to n. Then $\displaystyle\sum_{n=1}^{56}f(n)$ is equal to :-

\qn{6} If $\displaystyle\lim_{x\to0}\dfrac{ax-(e^{4x}-1)}{ax(e^{4x}-1)}$ exists and is equal to b, then the value of a $-$ 2b is \underline{\hspace{1cm}}.

\qn{7} If A and B are two sets such that A = \{1, 2, x\}, B = \{3, 4, y\} and \{1, 3\} $\times$ \{2, 4\} $\subseteq$ A $\times$ B then the value of x + y is :

\qn{8} The integral value of a for which the point ($-$2, a) lies in the interior of the triangle formed by the lines y = x, y = $-$x and 2x + 3y = 6.

\qn{9} If the lines $\dfrac{x-1}2=\dfrac{y+1}3=\dfrac{z-1}4$ and $\dfrac{x-3}1=\dfrac{y-k}2=\dfrac z1$ intersect, then 2k is equal to :

\qn{10} $96\cos\dfrac\pi{33}\cos\dfrac{2\pi}{33}\cos\dfrac{4\pi}{33}\cos\dfrac{8\pi}{33}\cos\dfrac{16\pi}{33}$ is equal to

\end{multicols}
\end{document}
